\section*{Metode}
\setcounter{section}{4}
\addcontentsline{toc}{section}{Metode}

\setcounter{subsection}{0}

Forsøket ble gjennomført ved å undersøke turbinens virkningsgrad for to ulike nålposisjoner. Pumpeturtallet ble først stilt inn til omtrent 2900 rpm. Deretter ble volumstrømmen justert til omtrent 170 L/min for nålposisjon 7 og 115 L/min for nålposisjon 3.

For hver nålposisjon ble belastningen på turbinakselen økt trinnvis ved å stramme bremsereimen. Ved hvert lastpunkt ble turbinens turtall, dreiemoment, volumstrøm og innløpstrykk lest av etter at målingene hadde stabilisert seg. Disse måleverdiene ble brukt til å beregne mekanisk effekt, hydraulisk effekt og virkningsgrad.

Forsøket ble først gjennomført for nålposisjon 7 og deretter gjentatt for nålposisjon 3. Resultatene ble til slutt sammenlignet ved å plotte virkningsgrad som funksjon av turbinturtall.
